\PassOptionsToPackage{table}{xcolor}
\documentclass[sigconf,review]{acmart}
\settopmatter{printacmref=false} % Removes citation information below abstract
\renewcommand\footnotetextcopyrightpermission[1]{} % removes footnote with conference information in first column
\pagestyle{plain} % removes running headers

\usepackage{pifont} % used for \todo symbol
\usepackage[normalem]{ulem} % in the preamble

\newcommand{\todo}[1]{\textcolor{blue}{\uline{#1}}}

\usepackage{hyperref}
\usepackage{float}
\usepackage{geometry}
\usepackage{array}
\usepackage{booktabs}
\usepackage{enumitem}
\usepackage{graphicx} % Required for inserting images
\usepackage{tikz}
\usepackage{pgfplots}

\usepackage{pgfplotstable}
\usepackage{tabularx}
\usetikzlibrary{shapes,arrows,positioning}
\usepackage{natbib}
\usepackage{chronology}
\usepackage{svg}
\usepackage{colortbl} % Add this line to your preamble
\usepackage{multirow}
\usepackage{makecell}


\makeatletter
\renewcommand\subsubsection{\@startsection{subsubsection}{3}{\z@}%
                                     {-3.25ex\@plus -1ex \@minus -.2ex}%
                                     {1.5ex \@plus .2ex}%
                                     {\normalfont\itshape}}

\makeatother





\def\BibTeX{{\rm B\kern-.05em{\sc i\kern-.025em b}\kern-.08emT\kern-.1667em\lower.7ex\hbox{E}\kern-.125emX}}

% Rights management information.
% This information is sent to you when you complete the rights form.
% the commands and values with those provided to you when you complete the rights form.
%
% These commands are for a PROCEEDINGS abstract or paper.

% \copyrightyear{2018}
% \acmYear{2018}
% \setcopyright{acmlicensed}
% \acmConference[Woodstock '18]{Woodstock '18: ACM Symposium on Neural Gaze Detection}{June 03--05, 2018}{Woodstock, NY}
% \acmBooktitle{Woodstock '18: ACM Symposium on Neural Gaze Detection, June 03--05, 2018, Woodstock, NY}
% \acmPrice{15.00}
% \acmDOI{10.1145/1122445.1122456}
% \acmISBN{978-1-4503-9999-9/18/06}

%
% These commands are for a JOURNAL article.
%\setcopyright{acmcopyright}
%\acmJournal{TOG}
%\acmYear{2018}\acmVolume{37}\acmNumber{4}\acmArticle{111}\acmMonth{8}
%\acmDOI{10.1145/1122445.1122456}

%
% Submission ID.
% Use this when submitting an article to a sponsored event. You'll receive a unique submission ID from the organizers
% of the event, and this ID should be used as the parameter to this command.
%\acmSubmissionID{123-A56-BU3}

%
% The majority of ACM publications use numbered citations and references. If you are preparing content for an event
% sponsored by ACM SIGGRAPH, you must use the "author year" style of citations and references. Uncommenting
% the next command will enable that style.
%\citestyle{acmauthoryear}

% The first command in your LaTeX source must be the \documentclass command.
%
% end of the preamble, start of the body of the document source.
\begin{document}

% These commands have SAMPLE values in them; it is your responsibility as an author to replace
  %
% The "title" command has an optional parameter, allowing the author to define a "short title" to be used in page headers.
\title{Bridging HPC and Cloud: Issues and Incentives}

%
% The "author" command and its associated commands are used to define the authors and their affiliations.
% Of note is the shared affiliation of the first two authors, and the "authornote" and "authornotemark" commands
% used to denote shared contribution to the research.

% \author{Ben Trovato}
% \authornote{Both authors contributed equally to this research.}
% \email{trovato@corporation.com}
% \orcid{1234-5678-9012}
% \author{G.K.M. Tobin}
% \authornotemark[1]
% \email{webmaster@marysville-ohio.com}
% \affiliation{%
%   \institution{Institute for Clarity in Documentation}
%   \streetaddress{P.O. Box 1212}
%   \city{Dublin}
%   \state{Ohio}
%   \postcode{43017-6221}
% }

\author{German Savchenko and Tiziano De Matteis}
\affiliation{%
  \institution{Vrije Universiteit Amsterdam, The Netherlands}
}
\email{g,savchenko@student.vu.nl, t.de.matteis@vu.nl}


%
% By default, the full list of authors will be used in the page headers. Often, this list is too long, and will overlap
% other information printed in the page headers. This command allows the author to define a more concise list
% of authors' names for this purpose.
%\renewcommand{\shortauthors}{Verdecchia, et al.}

%
% The abstract is a short summary of the work to be presented in the article.
\begin{abstract}

%\todo{The abstract of the study goes here. A good abstract shortly summarizes the paper by covering: (i) a presentation of the topic covered, (ii) the research followed, (iii) the main findings of the study, and (iv) the outcome of the investigation}

    The shift of High Performance Computing (\textit{HPC}) workloads from \textit{on-premise} to \textit{cloud} environments is a critical topic, driven by the demand for scalable, flexible, and cost-effective computing solutions. An overview of the challenges and opportunities associated with this transition are provided, focusing mostly on research done within the last 5 years. To gain a broader perspective on the current state of the field, we also consider older research that highlights the initial limitations faced in earlier stages of HPC-Cloud convergence and attempts to delineate a timeline of progress.
    Various aspects of \textit{HPC-Cloud convergence} are examined, as well as how it is affected by network performance, virtualization technologies, specialized hardware, containerization, cloud storage solutions, resource disaggregation, and emerging trends like serverless computing, edge computing, and AI integration. Key metrics identified include network latency, virtualization overhead, unpredictable performance, cost considerations, limited hardware options, and the lack of HPC-specific tools. The review highlights the importance of robust security measures, efficient data management and access, intelligent resource allocation and scheduling, and application adaptability and portability.
    We conclude that a hybrid approach, combining on-premise HPC with Cloud resources, offers the most practical and beneficial path for convergence. This hybrid approach enables cloud bursting, workflow management, and software resource disaggregation, leveraging the strengths of both environments while mitigating the risks of a complete cloud transition. The study identifies promising advancements in cloud networking, containerization technologies, specialized hardware, and cloud storage services, contributing to a future where HPC and cloud computing increasingly work together. This integration enhances the reproducibility of scientific research, broadens access to high-performance computing for more organizations and researchers, and drives new opportunities for scientific discovery and technological innovation.



\end{abstract}

%
% The code below is generated by the tool at http://dl.acm.org/ccs.cfm.
% Please copy and paste the code instead of the example below.
%
% \begin{CCSXML}
% <ccs2012>
%  <concept>
%   <concept_id>10010520.10010553.10010562</concept_id>
%   <concept_desc>Computer systems organization~Embedded systems</concept_desc>
%   <concept_significance>500</concept_significance>
%  </concept>
%  <concept>
%   <concept_id>10010520.10010575.10010755</concept_id>
%   <concept_desc>Computer systems organization~Redundancy</concept_desc>
%   <concept_significance>300</concept_significance>
%  </concept>
%  <concept>
%   <concept_id>10010520.10010553.10010554</concept_id>
%   <concept_desc>Computer systems organization~Robotics</concept_desc>
%   <concept_significance>100</concept_significance>
%  </concept>
%  <concept>
%   <concept_id>10003033.10003083.10003095</concept_id>
%   <concept_desc>Networks~Network reliability</concept_desc>
%   <concept_significance>100</concept_significance>
%  </concept>
% </ccs2012>
% \end{CCSXML}

% \ccsdesc[500]{Computer systems organization~Embedded systems}
% \ccsdesc[300]{Computer systems organization~Redundancy}
% \ccsdesc{Computer systems organization~Robotics}
% \ccsdesc[100]{Networks~Network reliability}

%
% Keywords. The author(s) should pick words that accurately describe the work being
% presented. Separate the keywords with commas.
\keywords{Systematic Literature Review, High Performance Computing (HPC), Cloud computing, HPC-Cloud Convergence, Hybrid Cloud, HPC-Cloud transition challenges}

%
% A "teaser" image appears between the author and affiliation information and the body
% of the document, and typically spans the page.

%
% This command processes the author and affiliation and title information and builds
% the first part of the formatted document.
\maketitle
% \tableofcontents
%\vspace{10pt}
\section{Introduction}
% \todo{Present the topic of the literature review, including background of research field, goal, methodology, and contribution of the paper (e.g., map of the state of the art, reusable classification approach, evaluation of results, discussion of results, and target audience)}

For a good part of the last decade the interest in transitioning from HPC on-premise infrastructures to an on-demand Cloud environments has been intensifying and pulling in numerous different directions. The research regarding the topic is continuing to grow and  both HPC and Cloud are distinctively improving on their core technologies, as well as growing in popularity. As we saw the exascale milestone was officially being achieved in mid 2022, and the yearly growth of both markets is high. Currently, most HPC platforms, both academic and industrial, are on-premise. However, a report by Hyperian Research states that while the on-premise HPC market is growing at 7\% a year until 2024, the HPC Cloud market will grow at 17\% a year \cite{silva_review_2024} in the Netherlands. In the U.S., holding the largest share of global market share in both HPC and Cloud, the prospects of annual growth are even higher (20\% for Cloud and 9\% for HPC), while the market shares are (for 2023) \$14.6 billion for HPC and \$216 billion for Cloud. This puts the process of convergence in perspective, suggesting that while HPC remains a critical and distinct technology, it is the one being integrated within the broader ecosystem of Cloud computing.



There are multiple areas of interest within the HPC-Cloud discussion and most of them seem to primarily branch from the performance inconsistencies when transitioning from one system to another, as well as the technical challenges in adapting from one system to another.

\noindent
\begin{enumerate}
\item \textbf{Containerization and Virtualization:} These pivotal technologies facilitate the transition, allowing containerized software to uniformly operate across diverse platforms, fostering the seamless exchange of technology between HPC and Cloud. The adaptability introduced by containerization and virtualization plays a crucial role in easing the transition process.
\item \textbf{Performance Evaluation:} The performance evaluation of distinct HPC-Cloud configurations and their convergence emerges as a critical realm of exploration. \textit{HPC-Cloud convergence} will be used throughout the text and indicates the intersection and co-working of HPC and Cloud paradigms. The inherent challenges associated with performance inconsistencies during the transition from one system to another are still unclear. Unraveling these intricacies is necessary for optimizing the performance of HPC applications in Cloud environments or in mixed environments, ensuring their efficacy and reliability. The heterogenization of large compute systems adds another layer of complexity to the evolving landscape of HPC and Cloud computing. Specialized hardware, such as CPUs, GPUs, Tensor Processing Units (TPUs), Field Programmable Gate Arrays (FPGAs), and other domain-specific hardware, reshape the traditional paradigms in computing and, thus, are hard to bridge altogether.

\item \textbf{Reproducibility of Scientific Results:} Ensuring the reproducibility of scientific results remains a fundamental requirement within the context of HPC-Cloud integration. Robust methodologies must be in place to guarantee that scientific experiments conducted in Cloud environments can be accurately reproduced on different platforms or at various points in time, upholding the integrity of scientific research. Data access transparency emerges as a paramount concern in an era where scientific data is increasingly stored and processed in the Cloud. Establishing mechanisms for transparent and efficient access to this data is imperative to sustain the fluidity of scientific workflows. The user experience and migration frameworks play a pivotal role in the successful adoption of Cloud technologies within the HPC community and vice-versa. Developing tools and frameworks that facilitate the seamless migration of HPC applications to the Cloud, coupled with an emphasis on providing an intuitive and user-friendly experience, are essential components for fostering widespread acceptance and utilization.

\item \textbf{Economic Considerations:} Navigating the economic landscape, the costs associated with transitioning to the Cloud in a dynamic market and evolving infrastructures come to the forefront. Organizations contemplating this shift must carefully consider the economic aspects of Cloud computing, focusing on cost-effectiveness and evaluating the return on investment. Despite the growing prominence of Cloud computing, there exists a notable gap in research concerning large-scale workloads of HPC in the Cloud.
\end{enumerate}

Overall, the transition from on-premise HPC to the Cloud is a multifaceted journey, encompassing a spectrum of high-stakes considerations. This review attempts to shed some additional light on what is the current estimation of the future for the co-existence and merging of Cloud and on-premise systems, focusing mainly on current challenges that both the stakeholders and the users encounter in the process of transitioning. Additionally, an overview of popular challenges of HPC in the Cloud (e.g. network bandwidth and latency, topology, multi-tenancy, virtualization, and filesystem issues, among others) is examined to understand whether they persist at this moment in time or have been solved by recent innovation.

% containerization, virtualization, and exchange of technology of the two fields
% performance of different HPC-Cloud configurations and convergence
% heteronization of hardware
% reproducibility of scientific results
% data access transparency
% costs in a constantly evolving market and infrastructures
% lack of studies on large scale workloads of hpc in Cloud
% sustainability issues
% user experience and migration frameworks

% “As heterogeneity increases, so does the risk of underutilizing expensive hardware resources if we preserve today’s rigid node configuration and reservation strategies.” ([Michelogiannakis et al., 2022, p. 291](zotero://select/library/items/NQ9MZGGT)) ([pdf](zotero://open-pdf/library/items/UWPKERIF?page=1))

\section{Related Work}
\label{sec:related_work}
% \todo{Section dedicated to the related work, i.e., similar / relevant studies, and how your study fits in the overall research picture. Usually, this section includes secondary studies which address the same / similar topic to the one addressed in your study.}

This systematic literature review builds upon a growing body of research exploring the convergence of HPC and cloud computing. Several existing studies have investigated various aspects of this transition, providing valuable insights into the challenges, opportunities, and future directions of this field.

\noindent
\textbf{Early Research on HPC in the Cloud}
Early research focused on identifying the challenges of migrating HPC applications to cloud environments. Studies like \cite{wang2010cloud, gupta2014challenges, marathe2013challenges} highlighted the limitations of cloud networks, virtualization overhead, and the unpredictable nature of shared cloud resources. These studies emphasized the need for specialized tools and infrastructure to support HPC workloads in the cloud.

\noindent
\textbf{Recent Research on HPC-Cloud Convergence}
More recent research has explored the development of solutions and technologies to address the challenges identified in earlier studies. For example, \cite{netto_hpc_2019} provides a comprehensive overview of the state of the art in HPC-Cloud convergence, highlighting the importance of hybrid cloud models, containerization technologies, and resource disaggregation.
Several studies have focused on specific aspects of HPC-Cloud convergence, such as: \\
\newline
\noindent
\textit{Network Performance}: \cite{alves_interference-aware_2020, giannoutakis_evaluation_2021} investigate the performance of cloud networks for HPC workloads, exploring the use of high-speed interconnects and optimized communication protocols.\\

\noindent
\textit{Virtualization Technologies}: \cite{marques_harnessing_2021} examines the use of containerization and operating-system-level virtualization to reduce overhead and improve performance in cloud-based HPC environments.\\

\noindent
\textit{Resource Disaggregation}: \cite{michelogiannakis_case_2022, copik_software_2024} explore the potential of resource disaggregation to optimize hardware utilization and enhance flexibility in HPC-Cloud systems.\\

\noindent
\textit{Cloud Storage for HPC}: \cite{gadban_analyzing_2022, guidi_10_2021} investigate the use of cloud storage services like Amazon S3 for HPC workloads, highlighting the benefits and challenges of this approach.\\

\noindent
\textit{AI and Machine Learning}: \cite{chahal_isesa_2022, michalewicz_attempt_2023, silva_review_2024} explore the use of AI and ML for optimizing HPC workloads, managing cloud resources, and enabling new scientific applications.\\

This review builds upon the existing literature by providing a comprehensive overview of the current state of the art in HPC-Cloud convergence, focusing on recent research (2021-2024) and highlighting emerging trends. The review also emphasizes the importance of a hybrid approach, combining on-premise HPC with cloud resources, as a practical and beneficial path for achieving convergence. Furthermore, the review explores the potential of AI and ML to address key challenges and drive innovation in the field.
By synthesizing and analyzing the findings of previous studies, this review aims to provide a valuable resource for researchers, practitioners, and policymakers interested in the future of HPC-Cloud convergence.

\section{Study Design}\label{sec:design}
    % Section dedicated to document your study was carried out
The following subsections provide an account of the initial search strategy (\autoref{sec:init_search}), selection criteria (\autoref{sec:sel_criteria}), data extraction process (\autoref{sec:data_extraction}), data synthesis approach (\autoref{sec:data_synth}), and steps taken to ensure the replicability of the study (\autoref{sec:replicability}).


\subsection{Research Goal}
\label{sec:research_goal}
    % \todo{What is the research goal of this study should be documented in this section. This can follow the GQM approach, as described in the following sentence, which NEEDS TO BE REPHRASED IN YOUR OWN WORDS if you want to include it in your study.}

    The topic of HPC in the Cloud has been studied for some time now, and challenges of infrastructure, performance, portability, and costs persisted for a good part of the last decade. For example, the lack of fast intra-node networks in the Cloud, often relying on TCP/IP and Ethernet, compared to the \textit{high-performance interconnects} used in HPC clusters (e.g., InfiniBand, Cray Slingshot, RDMA over Converged Ethernet (RoCE)), significantly impacted the performance of tightly coupled HPC applications \cite{wang2010cloud, gupta2014challenges}. These network bottlenecks also increased the cost of running long-running HPC applications in the Cloud due to unpredictable performance variations \cite{gupta2014challenges, marathe2013challenges, de_sensi_noise_2022}. Additionally, the shared nature of Cloud resources made it difficult to predict the performance of HPC applications, as the presence of other workloads on the same physical infrastructure could lead to unpredictable performance variations and resource contention \cite{gupta2014challenges, marathe2013challenges}. These factors, along with the cost of Cloud computing for HPC workloads, hindered the adoption of Cloud solutions for HPC in the early stages \cite{jackson2010performance, gupta2014challenges, marathe2013challenges}.

    The primary objective of this literature review, thus, is to gain insight into the current challenges and potential solutions associated with transitioning on-premise HPC applications to the Cloud or the convergence of the two. The study also entails defining exactly the main areas of interest at the moment in terms of the effort in bringing the two types of systems closer, whether the trend of \textit{Cloudification} is promising, and analyzing studies that focused on some very common HPC applications, like AI workloads, scientific simulations, and so on.  To underline the main objective of this review, we define the following contributions of this report.

    \begin{center}
        \begin{itemize}
            \item [\textbf{Con-1}] An update of the current, fast-changing, state of the art in HPC-Cloud interaction
            \item [\textbf{Con-2}] A qualitative advice for researchers and scientific programmers on what to expect from their large-scale applications when moving from on-premise to Cloud systems
            \item [\textbf{Con-3}] An overview of what are the promising convergence modalities between HPC and Cloud
            \item [\textbf{Con-4}] An update of the current trends, challenges, and expectations of the HPC-Cloud field for the end of the decade
        \end{itemize}
    \end{center}


\subsection{Research Questions}
\label{sec:questions}
    % \todo{In this subsection, the research questions underlying your study should be reported}
We want to define and cover the evolution of the HPC-Cloud transition and their merging areas, and investigate the latest advances and results from influential publications. Based on this premise, we define the following Research Questions:\\


    % \todo{to update based on feedback}
\begin{itemize}
    \item \textbf{RQ-1}: What are the key challenges in transitioning on-premise HPC applications to the Cloud?
    \item \textbf{RQ-2}: How did the transition challenges evolved in time?
    \item \textbf{RQ-3}: What are the main HPC-Cloud configurations, and how do they affect performance, cost, and energy? What are the trade-offs?
    \item \textbf{RQ-4}: What hardware and software design aspects of HPC applications hinder seamless Cloud integration, and how are these challenges being addressed?
\end{itemize}

    % \todo{to update based on feedback}



\subsection{Initial search}
\label{sec:init_search}
    % \todo{How was the initial search for literature carried out should be reported in this section}
    To adhere to the initially defined objectives, the study aimed to identify a comprehensive set of studies that explored the convergence of HPC and Cloud computing. This search was conducted using a combination of databases and keywords, carefully selected to capture the breadth and depth of the research area.
    The primary databases used were Scopus, IEEE Xplore, and dblp, chosen for their comprehensive coverage of computer science and engineering research. These databases offer a vast repository of peer-reviewed publications, ensuring the inclusion of high-quality and relevant studies in the review.
    The selection of keywords was equally important, as they served as the foundation for the initial search and guided the identification of relevant studies. The keywords were chosen to reflect the key themes of the review, encompassing both general and specific aspects of HPC-Cloud convergence.
    The initial search employed a combination of general keywords, such as "High-Performance Computing", "Cloud Computing", "HPC in the Cloud", and "HPC-Cloud Convergence", to capture the broad scope of the research area. These keywords were used to identify studies that addressed the overall concept of integrating HPC and Cloud computing.
    To further refine the search and identify studies that focused on specific aspects of convergence, additional keywords were incorporated. These included "Cloud Bursting", "Resource Disaggregation", "AI for HPC", "Machine Learning for HPC", and "Serverless Computing" \autoref{tab:keywords}. These keywords helped to identify studies that explored specific challenges, solutions, and emerging trends within the field of HPC-Cloud convergence \cite{gadban_analyzing_2022, michelogiannakis_case_2022, copik_software_2024, chahal_isesa_2022, michalewicz_attempt_2023, silva_review_2024, liang_scalable_2022, eberius_evaluating_2023, alves_interference-aware_2020, kolodziej_cloud-based_2019, kabir_uncertainty-aware_2022, reed_hpc_2023}.
    The initial search yielded numerous results from the past 5 years, reflecting growing interest and research activity in this dynamic field. However, there was almost no material offering a comprehensive, high-level overview of the convergence of the two paradigms (\cite{netto_hpc_2019}, \cite{liu_hpc_2022}). This initial set of results provided a starting point for the review, but refining the selection using the criteria described in the next section was essential. These criteria ensured that only the most relevant and high-quality studies were included, ultimately contributing to a comprehensive and insightful analysis of the current state of HPC-Cloud convergence.

\begin{table*}[ht!]
\centering
\small
\begin{tabular}{c|l|c|c}
\hline
\rowcolor{gray!20} \textbf{No.} & \textbf{Keywords} & \textbf{Number of Articles Found} & \textbf{Articles Used} \\ \hline
\hline
\textbf{First Iteration} \\ \hline
\rowcolor{white} 1 & HPC-Cloud & 89 & 15 \\ \hline
\rowcolor{gray!20} 2 & HPC AND cloud AND convergence & 38 & 8 \\ \hline
\hline
\textbf{Second Iteration} \\ \hline
\rowcolor{white} 3 & HPC AND Cloud AND Bursting & 10 & 3 \\ \hline
\rowcolor{gray!20} 4 & HPC AND Resource AND Disaggregation & 22 & 7 \\ \hline
\rowcolor{white} 5 & Serverless AND HPC & 53 & 7 \\ \hline
\rowcolor{gray!20} 6 & Reproducibility AND HPC AND Cloud & 41 & 2 \\ \hline
\rowcolor{white} 7 & Elastic AND Cloud AND HPC & 47 & 5 \\ \hline
\end{tabular}
\caption{Number of articles found and used for each keyword}
\label{tab:keywords}
\end{table*}





\subsection{Application of selection criteria}
\label{sec:sel_criteria}


    As previously mentioned, the topic of HPC-Cloud is a very prolific area of research at the moment, and we can observe numbers of pre-prints (on Scopus) Just for the current year in order of four to six hundred for the selected search keywords. Therefore, we found it necessary to define a relatively large number of strict selection criteria in order to be able to satisfy the expected contributions mentioned in \autoref{sec:research_goal}.

    \begin{enumerate}
        \item[{I}1-] A study that addresses specific software or hardware solutions in HPC-Cloud context
        \item[{I}2-] A study that is peer-reviewed
        \item[{I}3-] A study that has been published in the last 5 years time frame
        \item[{I}4-] A study that brings together HPC and Cloud - direct/indirect comparison
        \item[{I}5-] A study that presents clear empirical metrics
    \end{enumerate}

    Similarly, we define the exclusion criteria:

    \begin{enumerate}
        \item[{E}1-] A study that is not peer-reviewed (i.e. not published on any reputable journal platform.)
        \item[{E}2-] A study that is outdated. In the scope of this study, the main points should be covered by literature in the range 2021-2024, while all the others can be supported by older literature for the sake of providing context and background.
        \item[{E}3-] A study with less than 10-20 citations. This aids in filtering out a large amount of non-primary papers, and to avoid possible threats to conclusion validity (\autoref{sec:threatsvalidity}).
    \end{enumerate}



\subsection{Snowballing}
\label{sec:snowballing}
    % \todo{If a snowballing approach was used for the literature review, its details should be documented in this subsection}

    In an attempt to gather a broad overview of the topic and relevant papers, a partial snowballing approach was adopted. Without a seed paper, the initial papers of interest helped define the search path. This approach was beneficial for understanding the overall breadth of the topic, the average age of the referenced articles, and identifying key figures, institutions, or journals that appeared repeatedly. It quickly helped gather material for the reference manager and prioritize essential sources. As mentioned in \autoref{sec:sel_criteria}, this method also ensured the inclusion of articles within the last 5 years.


\subsection{Data Extraction}
\label{sec:data_extraction}
    % \todo{Report in this section the data extraction followed to gather the data for the study (e.g., what process did you followed to gather the data in the companion data extraction spreadsheet?)}
    After an initial substantial body of research papers was gathered, the next step was to start to sort out the useful papers and to discard the less useful ones. The initial processing was done by reading the introduction, similar work, and conclusions of the selected papers. Furthermore, the idea was to combine the process of categorizing the papers as useful and unuseful with the process of forming classes of work. These classes would naturally form as the reading progressed, allowing us to form a restricted number of macro categories that the works would discuss. These categories include past challenges, infrastructure and technology improvements, deployment models (i.e. hybrid Cloud, HPCaaS, etc.)(\autoref{tab:hpc_cloud_categories}), economic aspects, portability, and will be discussed in greater detail in the following sections of this report.

    The process of data extraction was facilitated by a literature matrix, where comments and summaries were added, and the final inclusion of the paper into the useful stack could be signed. For other recaps, notes, topic tags, etc., my reference manager of choice, Zotero, was used. Here, the intuitive note tool allowed me to have all of the recap material smoothly gathered in one project directory, which allowed for an easy recollection of information throughout the writing of this review. Furthermore, the tags feature allowed me to easily categorize the papers based on the class of work and, subsequently, visualize all the relevant material for a specific class in one list.

\begin{table*}[ht]
    \centering
    \renewcommand{\arraystretch}{1} % Reduce vertical spacing
    \resizebox{0.8\textwidth}{!}{ % Scale the table to 80% of text width, adjust height automatically
    \small
    \begin{tabularx}{\textwidth}{>{\raggedright\arraybackslash}p{4.5cm} >{\raggedright\arraybackslash}X >{\raggedright\arraybackslash}X}
    \toprule
    \textbf{Category} & \textbf{Challenges} & \textbf{Solutions} \\
    \midrule
    \textbf{Technological Aspects and Infrastructure} &
    \begin{itemize}[label={\rule[0.1ex]{0.8ex}{0.8ex}}]
    \item \textbf{Network Latency} Latency in cloud networks hinders performance of tightly coupled HPC applications.
    \item \textbf{Virtualization Overhead} Virtualization introduces overhead, impacting HPC performance.
    \item \textbf{Heterogeneity of Hardware}  Diverse accelerators and memory technologies create challenges for efficient resource utilization.
    \item \textbf{Data Access Transparency} Ensuring seamless data access across different storage systems is crucial.
    \item \textbf{User Experience and Migration Frameworks}  Providing a user-friendly interface and simplifying migration processes for HPC applications.
    \end{itemize} &
    \begin{itemize}[label={\rule[0.1ex]{0.8ex}{0.8ex}}]
    \item High-speed interconnects (InfiniBand, RDMA), optimized protocols (MPI, RDMA), new technologies (Slingshot, CXL, Optical)
    \item Containerization (Docker, Apptainer), OS-level virtualization, Bare-metal servers\newline
    \item Cloud providers offering a wider range of specialized hardware, bare-metal servers with custom configurations
    \item Libraries like S3Embedded, standardized protocols (S3 API), parallel file systems (Lustre, DAOS)
    \item User-centric design of converged systems, containerization for portability, automated tools for deployment and management
    \end{itemize} \\
    \midrule
    \textbf{Performance and Configuration} &
    \begin{itemize}[label={\rule[0.1ex]{0.8ex}{0.8ex}}]
    \item \textbf{Performance of HPC-Cloud Configurations} Optimizing HPC application performance across different cloud configurations.
    \item \textbf{Scaling HPC Workloads} Addressing challenges in running large-scale HPC workloads in cloud environments.
    \end{itemize} &
    \begin{itemize}[label={\rule[0.1ex]{0.8ex}{0.8ex}}]
    \item Advanced scheduling and allocation algorithms, dynamic resource management (DMR API, PMIx), cloud bursting\newline
    \item High-performance cloud networks, elastic jobs, efficient resource management systems
    \end{itemize} \\
    \midrule
    \textbf{Reproducibility and Transparency} &
    \begin{itemize}[label={\rule[0.1ex]{0.8ex}{0.8ex}}]
    \item \textbf{Reproducibility of Scientific Results} Ensuring consistent and reproducible execution of HPC applications.
    \item \textbf{Sustainability} Balancing computational demands with environmental concerns.
    \end{itemize} &
    \begin{itemize}[label={\rule[0.1ex]{0.8ex}{0.8ex}}]
    \item Containerization, standardized formats and APIs (OCI), version control systems\newline
    \item Energy-efficient hardware, workload optimization, spatio-temporal workload shifting
    \end{itemize} \\
    \midrule
    \textbf{Economic and Market Dynamics} &
    \begin{itemize}[label={\rule[0.1ex]{0.8ex}{0.8ex}}]
    \item \textbf{Cost Efficiency} Balancing performance with cost considerations in cloud-based HPC solutions.
    \item \textbf{Market Evolution} Understanding the evolving cost landscape and infrastructure changes in the cloud computing market.
    \end{itemize} &
    \begin{itemize}[label={\rule[0.1ex]{0.8ex}{0.8ex}}]
    \item Pay-as-you-go pricing models, resource disaggregation, automated optimization systems
    \item Keeping abreast of industry trends, engaging with cloud providers, exploring new technologies
    \end{itemize} \\
    \bottomrule
    \end{tabularx}
    }
    \caption{Categories of Consideration in High-Performance Computing in the Cloud}
    \label{tab:hpc_cloud_categories}

\end{table*}



% \begin{table*}[ht]
%     \centering
%     \begin{tabular}{>{\raggedright}p{3.5cm} p{10cm}}
%     \toprule
%     \textbf{Category} & \textbf{Considerations} \\
%     \midrule
%     \textbf{Technological Aspects and Infrastructure} &
%        \begin{itemize}[label={\rule[0.2ex]{1.0ex}{1.0ex}}]
%           \item Containerization, virtualization, and exchange of technology between HPC and Cloud.
%           \item Heteronization of hardware.
%           \item Data access transparency.
%           \item User experience and migration frameworks.
%        \end{itemize} \\
%     \midrule
%     \textbf{Performance and Configuration} &
%        \begin{itemize}[label={\rule[0.2ex]{1.0ex}{1.0ex}}]
%           \item Performance of different HPC-Cloud configurations and convergence.
%           \item Lack of studies on large scale workloads of HPC in the Cloud.
%        \end{itemize} \\
%     \midrule
%     \textbf{Reproducibility and Transparency} &
%        \begin{itemize}[label={\rule[0.2ex]{1.0ex}{1.0ex}}]
%           \item Reproducibility of scientific results.
%           \item Sustainability issues.
%        \end{itemize} \\
%     \midrule
%     \textbf{Economic and Market Dynamics} &
%        \begin{itemize}[label={\rule[0.2ex]{1.0ex}{1.0ex}}]
%           \item Costs in a constantly evolving market and infrastructures.
%        \end{itemize} \\
%     \bottomrule
%     \end{tabular}
%     \caption{Categories of Consideration in High-Performance Computing in the Cloud}
%     \label{tab:hpc_cloud_categories}
%     \end{table*}




\subsection{Data Synthesis}
\label{sec:data_synth}
    % \todo{What approach did you followed to carry out the data synthesis process and summarize the data extracted from the primary studies?}
    As briefly mentioned in \autoref{sec:data_extraction}, most of the data synthesis was done using the reference manager software of choice, Zotero. After an initial screening of the papers and the addition of basic notes over the introduction and conclusion sections, the next phase would be to include a second set of lower-level summaries of the, first, some of the most important papers covering all classes of work, or covering one class of work in depth and with high grade of influence and quality. This would allow us to generate the first synthesized material, and to place it into the correct sections of classes of work. The classes of work, initially defined, were meant to be partly malleable and to allow room for updates, addition, removal, and modification in later stages, as the review was being synthesized.

    After defining the first review parts based on the most influential primary papers, the smaller and more specific papers would fill further details of the review.


\subsection{Study Replicability}
\label{sec:replicability}


    The reproducibility of this survey is made possible by making the used matrices in the data extraction phase accessible to the public \cite{literature-matrix}. The provided table includes the majority of the papers and their contents sorted in the respective macro categories mentioned in \autoref{sec:data_extraction}.

\subsection{Threats to Validity}
\label{sec:threatsvalidity}
% \todo{In this section, the most prominent threats to validity of your study should be discussed (the following categories are just an example suggestion). They have been moved with study design, as one should consider them before study execution.}

In the exploration of HPC in the Cloud one significant threat to validity is the potential presence of "implementation bias." This threat emerges when studies predominantly focus on successful implementations of HPC in the Cloud, possibly overlooking challenges, failures, or suboptimal instances.

For instance, researchers and industry professionals might be more inclined to publish and share success stories of HPC applications in the Cloud due to their positive outcomes. As a result, the literature may disproportionately highlight favorable scenarios, potentially giving a skewed perception of the actual challenges and limitations associated with implementing HPC solutions in cloud environments.

To address this threat, it is crucial for the literature review to actively seek out and include studies that discuss unsuccessful or challenging implementations of HPC in the Cloud. Incorporating a diverse range of perspectives and outcomes ensures a more comprehensive and realistic understanding of the successes and limitations within this dynamic intersection of High-Performance Computing and Cloud technology.

\subsubsection{External Validity}
In the context of external validity for HPC in the Cloud, one potential concern is the generalizability of findings across different cloud service providers and HPC architectures. Variability in infrastructure, performance capabilities, and resource allocation methods among different cloud platforms may impact the external validity of research findings. To enhance external validity, it is essential to consider a diverse range of cloud environments and HPC configurations in the literature review.

\subsubsection{Internal Validity}
Internal validity is crucial for ensuring the accuracy and reliability of findings within the realm of HPC in the Cloud. Potential threats may arise from variations in experiment design, such as differences in workload characteristics, benchmarking methodologies, and measurement tools. Researchers should carefully evaluate and account for these factors to maintain internal validity, providing a solid foundation for drawing accurate conclusions about the performance and efficiency of HPC applications in cloud environments.

\subsubsection{Construct Validity}
Construct validity in the study of HPC in the Cloud involves ensuring that the chosen measures and metrics accurately capture the intended aspects of performance, scalability, and efficiency. For example, many studies refer to large-scale applications while testing on small-scale infrastructures to limit costs, which is not contributing effectively to this literature review. Challenges may arise from the use of diverse performance metrics across studies, making direct comparisons difficult. Therefore, ensuring a standardized set of metrics and benchmarks relevant to HPC workloads in the Cloud can enhance construct validity, allowing for meaningful comparisons and a more cohesive understanding of the field.

\subsubsection{Conclusion Validity}
In the context of HPC in the Cloud pertains to the accuracy and appropriateness of inferences drawn from the reviewed literature. Potential threats include incomplete reporting of results, selective publication, or biased interpretation of findings. This is an important factor to consider given that the contributions listed in \autoref{sec:research_goal} might be compromised. To enhance conclusion validity, it is crucial to critically evaluate the methodologies and interpretations presented in the literature, ensuring that the conclusions drawn accurately reflect the nuances and complexities of implementing HPC in diverse cloud computing environments.

\section{Results}
\label{sec:results}
    % \todo{In this section, the results of the literature review should be reported. This section can be split up according to the specific focus of the literature review, e.g. Publication Trends, Research Focus, and Potential for Industrial Adoption (see Running example)}

\subsection{Early Challenges and Limitations}
\label{sec:earlychallenges}

    In the early days of Cloud computing, several factors made it difficult to port HPC applications to the Cloud. These challenges stemmed from the fundamental differences between traditional HPC environments and Cloud platforms, which were not initially designed to meet the demanding requirements of HPC workloads.
    The landscape has since evolved rapidly. Advancements in hardware and software, networking, containerization technologies, storage systems, and even automated integration of AI into processes have narrowed the performance gap between Cloud and HPC environments. This improvement is evident in the analyzed literature and the volume of new releases related to this topic (\autoref{tab:hpc_cloud_categories}).
    However, early challenges and limitations continue to influence the adoption and optimization of HPC applications in the Cloud. These challenges include: \\

    \noindent
    \textbf{Network Latency:} Cloud networks, often relying on TCP/IP and Ethernet, were designed for general-purpose applications, introducing significant latency compared to the high-speed interconnects used in HPC clusters (e.g., InfiniBand). This latency severely impacted the performance of tightly coupled HPC applications, where frequent communication between processes was crucial. Even with 10 Gigabit Ethernet, Cloud networks struggled to provide the low latency and high bandwidth needed for efficient HPC communication \cite{wang2010cloud}. This network bottleneck significantly increased the cost of running long-running HPC applications in the Cloud due to unpredictable performance variations \cite{gupta2014challenges, marathe2013challenges}. \\


    \noindent
    \textbf{Virtualization Overhead:} Virtualization technologies, used to create isolated environments for applications in the Cloud, introduced overhead in terms of CPU cycles, memory access, and network communication. This overhead further degraded the performance of HPC applications, especially those with high communication demands. The performance penalty associated with virtualization, particularly for HPC applications, was highlighted \cite{jackson2010performance}. \textit{Hypervisor-based virtualization}, a technology supporting multiple guest VMs running on a server and sharing resources, introduced significant overhead, impacting memory, I/O, and interconnect performance \cite{babu2014hypervisor, yunge2011cloud}. \\

    \noindent
    \textbf{Unpredictability:} The shared nature of Cloud resources made it difficult to predict the performance of HPC applications. The presence of other workloads on the same physical infrastructure could lead to unpredictable performance variations and resource contention. Predicting performance in a shared Cloud environment presented challenges \cite{gupta2014challenges, marathe2013challenges}. The presence of other workloads could significantly impact the performance of HPC applications, making it difficult to estimate execution times and costs.

    \noindent
    \textbf{Cost:} Cloud providers typically charge for compute and bandwidth usage, which could quickly add up for HPC workloads that generate large amounts of data. This made it challenging to estimate the cost of running HPC applications in the Cloud and hindered their adoption. The cost concerns associated with Cloud computing for HPC were acknowledged \cite{jackson2010performance, gupta2014challenges, marathe2013challenges}. The pay-as-you-go model, while flexible, could lead to high costs for resource-intensive HPC applications. \\

    \noindent
    \textbf{Limited Hardware Options:} Cloud providers offered a limited range of hardware options compared to HPC clusters, which often featured specialized processors and accelerators. This restricted the ability to run HPC applications that required specific hardware configurations. The limited hardware choices available in the Cloud compared to dedicated HPC systems were pointed out \cite{gupta2014challenges, marathe2013challenges}. This lack of specialized hardware options hindered the adoption of Cloud computing for certain HPC applications. \\

    \noindent
    \textbf{Lack of HPC-Specific Tools:} Cloud platforms lacked the tools and support needed for managing and deploying HPC applications. This included specialized schedulers, monitoring tools, and libraries for parallel I/O. The need for HPC-specific tools in the Cloud was emphasized \cite{yelick2011challenges, gupta2014challenges}. The lack of mature solutions for managing and deploying HPC applications hindered their adoption. \\

    While the analyzed literature is still rooted in these initial challenges, advancements in Cloud computing have addressed many of these issues. However, they continue to influence the design, optimization, and deployment of HPC applications in the Cloud. By understanding these early challenges, researchers and practitioners can develop strategies to overcome them and leverage the full potential of Cloud computing for HPC workloads.





    % ################################################################################
\subsection{Technological Aspects and Infrastructure}
\label{sec:techinfrastructure}

    % Containerization, virtualization, and exchange of technology between HPC and Cloud.
    % Heteronization of hardware.
    % Data access transparency.
    % User experience and migration frameworks.
    In order to better understand the interaction between the two paradigms we explore the relevant technologies and the exchange of these between HPC and Cloud. The underlying hardware, data management, and a general user experience in both environments are to be investigated too.
    We analyze, on the hardware level, network performance and specialized hardware (heterogeneity). Latency has been a great bottleneck for HPC workloads in the Cloud but this changed with the implementation of new \textit{interconnect technologies} inside Cloud server infrastructures, leading to a number of studies demonstrate near-bare-metal performance for some selection of workloads. Diversification of hardware within the Cloud servers has been limited by the lack of high speed interconnects \cite{alves_interference-aware_2020}\cite{michelogiannakis_case_2022}.
    On the software side, we analyze containerization, job elasticity, cloud storage and its usage in HPC, resource disaggregation, HPC specific tools in Cloud environments.





    % \subsubsection{Performance and scalability}
    % The convergence of HPC and Cloud computing is driven by the increasing demand for HPC resources and the need for flexible and scalable solutions. The reviewed literature highlights the significant progress made in bridging the performance gap between Cloud and HPC environments.

% \subsection{Hardware}
    \subsubsection{Hardware - Network Performance:} Historically, network latency and bandwidth limitations have been significant barriers to running HPC applications in the Cloud. However, advancements in Cloud networking technologies, including high-speed interconnects and optimized communication protocols such as \textit{InfiniBand} and \textit{RDMA} (Remote Direct Memory Access), have significantly narrowed the performance gap between Cloud and HPC. Advancements in Cloud networking technologies, including high-speed interconnects and optimized communication protocols, have significantly improved the performance of HPC applications in the Cloud. Studies like \cite{guidi_10_2021} and \cite{fernandez_evaluation_2022} demonstrate that modern Cloud platforms can achieve performance levels comparable to traditional HPC systems for specific workloads. The improvements in Cloud infrastructure, such as high-speed memory systems and interconnects, have significantly narrowed the performance gap. The studies compare the performance of the Cori Cray XC40 HPC system at NERSC with AWS EC2 instances optimized for memory-intensive and compute-intensive workloads. The results show that modern Cloud instances demonstrate significantly improved network performance, with AWS R5 achieving comparable or better inter-node communication performance than HPC systems. Additionally, optimized communication protocols and network configurations can further enhance the performance of Cloud-based HPC applications, making them suitable for a wide range of workloads. \cite{fernandez_evaluation_2022} analyzes the performance of parallel solvers and MPI communications using Infrastructure as a Service (IaaS) from public Cloud providers. The study utilizes the \textit{HPCG benchmark} and \textit{OSU microbenchmark} suite to assess the performance of HPC applications on IaaS from \textit{AWS}, \textit{Google Cloud Platform} (GCP), \textit{Azure}, \textit{Oracle Cloud Infrastructure}, and \textit{Packet}. The research highlights the significance of network bandwidth for scalability, demonstrating that clusters with 100 Gbps bandwidth maintain high scalability even at 4,500 cores. The study also emphasizes the importance of network latency and jitter for tightly coupled parallel solvers, showing that low-latency networks are essential for achieving high performance in Cloud-based HPC environments.

    While InfiniBand and RDMA continue to be key technologies (and implemented in Cloud environments \cite{maliszewski2020performance}), new advancements are emerging:
    \textit{Slingshot}: Cray's Slingshot interconnect, designed for exascale systems like Frontier and LUMI, offers high bandwidth and low latency, enabling efficient communication between nodes. It's adoption in Cloud, however, seems to be very limited at the moment.
    \textit{CXL (Compute Express Link)}: CXL is a new interconnect standard designed to provide high-speed communication between CPUs, GPUs, and other accelerators. It promises to improve the performance and flexibility of disaggregated HPC systems \cite{shan_hybrid_2022}. Additionally, the technology seems to be gaining popularity in Cloud environments \cite{10247726}.
   \textit{ Optical Interconnects}: Optical interconnects are gaining traction in HPC and Cloud environments due to their ability to achieve higher bandwidth (up to 800 Gbps) and lower latency compared to electrical interconnects \cite{michelogiannakis_case_2022, guo_understanding_2023, michelogiannakis_efficient_2023}. These interconnects are also being adopted in Cloud environments \cite{9793366}.

    % \subsubsection{Hardware - Storage}
    % \label{subsec:storage}

    \subsubsection{Hardware - Specialized Hardware}
    The early Cloud platforms offered limited hardware options, primarily focusing on general-purpose CPUs. However, the increasing demand for HPC workloads, particularly those involving machine learning and artificial intelligence, has driven Cloud providers to offer a wider range of specialized hardware.
    \textit{GPUs}: GPUs, initially designed for graphics rendering, have become powerful accelerators for computationally intensive tasks, especially in machine learning and scientific computing. Cloud providers like AWS, Azure, and GCP now offer a variety of GPU instances, allowing users to leverage the parallel processing power of GPUs for their HPC applications \cite{giannoutakis_evaluation_2021, balis_scientific_2021, malla_hpc_2019, fernandez_evaluation_2022}.
    \textit{TPUs}: TPUs are custom-designed ASICs (Application-Specific Integrated Circuits) specifically optimized for machine learning workloads. Google Cloud Platform offers TPU instances, providing significant performance gains for training and inference of large-scale machine learning models \cite{pena-monferrer_hpc-cloud_2021}.
    \textit{FPGAs}: FPGAs are reconfigurable hardware devices that can be programmed to perform specific tasks. Cloud providers like AWS and Xilinx offer FPGA instances, allowing users to customize hardware for specific applications and achieve high performance for specialized workloads \cite{giannoutakis_evaluation_2021}.
    ARM Processors: ARM processors are gaining traction in the Cloud, with AWS's Graviton processors demonstrating significant performance improvements for certain workloads. The adoption of ARM in HPC is still in its early stages, but it holds promise for future exascale systems \cite{fernandez_evaluation_2022, nowicki_performance_2021}.

    The emergence of Bare-Metal servers allows for more direct access to the hardware, enabling users to leverage specialized hardware configurations for their HPC workloads without incurring in virtualization overheads \cite{guidi_10_2021, cho_possibility_2019,kumar_containerized_2022, kuity_investigating_2023}.


% \subsection{Software}
    \subsubsection{Software - Containerization:} Containerization technologies have revolutionized application deployment and portability. Containers encapsulate applications and their dependencies, allowing them to run consistently across different environments, including HPC clusters and Cloud platforms. Docker is a popular open-source containerization platform that provides a lightweight and portable way to package and run applications. Docker containers are isolated environments that include all the necessary dependencies for an application to run, making it easy to deploy and manage applications across different systems \cite{beserra_performance_2015, cho_possibility_2019, kuity_investigating_2023}. Apptainer (former Singularity), on the other hand, is a containerization platform specifically designed for HPC environments. It provides native support for MPI, InfiniBand, and GPUs, making it ideal for running high-performance applications in containers. Apptainer also offers features for security and reproducibility, ensuring that applications run consistently across different systems \cite{roy_mashup_2022, petrosyan_serverless_2022, elisseev_multiscale_2022, zhou_containerisation_2023, ruhela_characterizing_2021}. Other contestants like CharlieCloud and Podman have also emerged as promising solutions for deploying and managing HPC applications in Cloud environments. Sarus, similarly to Apptainer, is also a container engine specifically designed for HPC systems which combines container portability with native HPC performance, offering features like direct access to compute and I/O devices and integration with batch resource management systems.  \cite{copik_rfaas_2023, benedicic_sarus_2019}

    Studies like \cite{ruhela_characterizing_2021} and \cite{kuity_investigating_2023} demonstrate that containerized HPC environments can achieve near-native performance for most subcomponents, making them suitable for HPC workloads. The study evaluates the performance of container technologies on up to 6,144 nodes with 344,064 processes, using benchmarks like Intel MPI Benchmarks (IMB) and the MILC (lattice quantum chromodynamics) and VPIC (Vector Particle-In-Cell) applications. The results indicate minimal overheads for containerized runs compared to bare metal. In some cases, the network performance is even higher than native performance, possibly due to lower scheduling and process management overhead. These technologies also offer multiple degree of encapsulation, where the user is provided with the choice of trade-off between portability and ease of use against performance and security \cite{kuity_investigating_2023, cho_possibility_2019, ruhela_characterizing_2021}.

    \citeauthor{zhou_containerisation_2023} provides a comprehensive survey of containerization technologies and orchestration strategies for HPC systems. The study highlights the differences between containerization in Cloud and HPC environments, particularly focusing on security and performance. The paper concludes that while containerization offers significant benefits for HPC, such as improved software management and reproducibility, challenges remain in security and performance. Future research should focus on developing container technologies that better integrate with HPC environments, addressing compatibility, security, and performance trade-offs.


    \subsubsection{Software - Elasticity}
    The concept of elastic jobs, which can dynamically resize their process count at runtime, has emerged as a promising approach to improve HPC system utilization and performance. The concept is very common in Cloud computing but not in HPC. The technologies that make it possible are the following:

    \textit{DMR API}: The DMR (Dynamic Resource Management) API is a standard interface for managing elastic jobs in HPC systems. It allows applications to request changes in their resource allocation, enabling them to adapt to varying system conditions \cite{iserte_dmr_2018}.
    \textit{PMIx}: PMIx (Process Management Interface for Exascale) is a library that provides a standardized interface for managing processes in HPC systems. It supports features like dynamic process creation, termination, and communication, which are essential for implementing elastic jobs \cite{castain_pmix_2018, huber_pmix_2022}.
    Elasticity has the potential to significantly improve HPC system utilization and performance by allowing applications to adapt to changing resource availability.
    \textit{Flux}: Flux is a next-generation workload management framework designed for supercomputers, HPC clusters, Cloud servers, and laptops. It features a modular design and combines hierarchical resource management with graph-based scheduling to improve the performance, portability, flexibility, and manageability of scientific workflows \cite{ahn_flux_2020}.


    \subsubsection{Software - Cloud Storage for HPC}
    Cloud storage services have become a viable alternative to traditional HPC storage systems, offering advantages in terms of scalability, cost, and accessibility. However, the performance of Cloud storage in HPC environments still needs to be carefully evaluated, as it can introduce overhead and limitations compared to HPC-specific storage solutions.

    Services like \textit{Amazon S3}, \textit{Google Cloud Storage}, and \textit{Azure Blob Storage} have become increasingly popular for HPC workloads, offering high scalability, cost-effectiveness, and ease of access. The use of Cloud storage for HPC workloads has been discussed, \cite{guidi_10_2021, guidi2021evaluation} and the performance of S3 in HPC environments has been further investigated, highlighting potential benefits and challenges \cite{gadban2021s3embedded, gadban2020http}.

    \textit{Amazon S3}: Amazon S3 is a highly scalable and durable object storage service that provides a RESTful API for accessing data. It offers features like data replication, encryption, and access control, making it suitable for storing and managing large data sets \cite{gadban2021s3embedded, guidi2021evaluation}. The S3 API remains the de facto standard for accessing Cloud storage. Cloud providers and open-source solutions offer S3-compatible storage services, making it a versatile option for HPC applications. The S3Embedded library, developed in this work, demonstrates the potential for seamless integration of Cloud storage into HPC systems. It offers a lightweight and efficient way to leverage existing HPC storage infrastructure for Cloud-based applications \cite{gadban2021s3embedded}.
    \textit{DAOS}: DAOS (Distributed Application Object Storage) is an open-source, parallel file system designed for high-performance storage operations. It is optimized for persistent memory technologies like Optane DIMMs and NVMe-accessed SSDs, offering low latency and high throughput \cite{wu2019daos, lofstead2016daos}.
    \textit{Google Cloud Storage}: Google Cloud Storage is a similar object storage service offered by Google Cloud Platform. It provides a RESTful API for accessing data and offers features like data replication, encryption, and access control \cite{gadban_investigating_2020, guidi_10_2021}.
    \textit{Azure Blob Storage}: Azure Blob Storage is Microsoft Azure’s object storage service. It provides a RESTful API for accessing data and offers features like data replication, encryption, and access control \cite{guidi_10_2021, endo_consideration_2022}.


    \subsubsection{Software - Resource disaggregation}

    The literature has shown that resource disaggregation is a relatively new field of research. The first paper that we found was published in 2013 \cite{barroso_resource_2013}. Since then, the number of publications has been increasing steadily. This indicates that resource disaggregation is a field of growing interest \cite{michelogiannakis_efficient_2023, michelogiannakis_case_2022, copik_software_2024}.

    The research focus of the papers that we found is mainly on the development of new algorithms and methods for resource disaggregation. The main outcomes found are better usage of racks, fewer network interfaces (NICs), fewer memory modules, and improved performance overall. Additionally, many studies focus on applying resource disaggregation with AI workloads. suggesting that machine learning in Cloud environments is a field of growing interest \cite{michelogiannakis_case_2022, michelogiannakis_efficient_2023, copik_software_2024}.

    The potential for industrial adoption of resource disaggregation, however, is not yet clear. The literature review has shown that there are some promising results, including improved power efficiency, but that there are also many open questions. For example, it is not yet clear how well resource disaggregation works in practice on a large scale (i.e. some experiments were executed on a set of 25 CPUs and 24 GPUs), and what the costs and benefits are. Further study is required for advancing the field forward.




    \subsubsection{Software - HPC-Specific Tools}
    Cloud providers have recognized the growing demand for HPC workloads and have developed tools and services specifically designed for managing and deploying HPC applications in the Cloud.
    \textit{AWS ParallelCluster}: AWS ParallelCluster is a tool that automates the deployment of HPC clusters on AWS. It simplifies the process of setting up a cluster with the necessary software and configurations, making it easier for HPC users to leverage AWS resources \cite{guidi2021evaluation}.
    \textit{Azure CycleCloud}: Azure CycleCloud is a similar tool offered by Microsoft Azure. It provides a user-friendly interface for deploying and managing HPC clusters on Azure, simplifying the process for HPC users \cite{guidi2021evaluation}.
    \textit{GCP Slurm Cluster}: Google Cloud Platform offers a tool for deploying Slurm clusters on Google Compute Engine. This tool simplifies the process of setting up a Slurm cluster on GCP, making it easier for HPC users to leverage GCP resources \cite{guidi_10_2021, guidi2021evaluation}.

    These tools have significantly reduced the barrier to entry for HPC users in the Cloud, making it easier to access and manage HPC resources.



    % IaaS (Infrastructure as a Service) provides on-demand remote access to client-specified compute environments, including both virtualized setups and "bare metal" setups for HPC, with CADD tools deployable as Cloud images or installed post-deployment;
    % PaaS (Platform as a Service) facilitates software operation and scaling in multi-user, multi-project environments; SaaS (Software as a Service) operates remotely, scaling across multiple nodes in an IaaS environment and presented typically as a web interface, widely adopted in various industries;
    %FaaS (Function as a Service), a SaaS component, is invoked without context, integrating VS-associated services into larger workflows
    %HaaS (Hardware as a Service)

    %([Olğaç et al., 2019, p. 270](zotero://select/library/items/PKURBRPY)).

     % ################################################################################
\subsection{Performance, Configuration, and Deployment Models}
\label{sec:perfconfiguration}

        % Performance of different HPC-Cloud configurations and convergence.
        % Lack of studies on large scale workloads of HPC in the Cloud.

        The interaction between an HPC system and a Cloud system exists on a spectrum. A system may lean more towards HPC with minimal integration of Cloud technologies, or vice versa, while others may strike a balance between the two.  Ideally, tasks on the \textit{mixed system} are executed with near-native performance, achieving efficiency as if running on the system specifically designed for them. The convergence, thus, can be categorized into three main approaches (\autoref{tab:hpc_cloud_approaches}): \textit{HPC in the Cloud}, which focuses on moving HPC applications to Cloud environments; \textit{HPC plus Cloud}, where Clouds are used to complement HPC resources \cite{endo_consideration_2022, balis_scientific_2021, XaaS-torsten}; and \textit{HPC as a Service}, which exposes HPC resources via Cloud services \cite{netto_hpc_2018, liu_hpc_2022}. A mixed system (converged HPC-Cloud) can handle tasks that were previously exclusive to individual non-converged systems, enabling a unified HPC and Cloud system to seamlessly execute both HPC and Cloud applications \cite{gadban_analyzing_2022}. A notable example is XaaS, a unified architecture leveraging performance-portable containers, which offers the flexibility and efficient resource utilization of serverless computing while accommodating long-running, performance-sensitive HPC workloads \cite{XaaS-torsten}.

        % Clouds provide rapid provisioning of resources and a significant reduction in operating costs related to energy, software licenses, and hardware obsolescence. With the recent increase in the number of cores in modern physical machines, HPC applications can scale up to many cores even when fully allocated to a single physical machine. Consequently, Clouds can offer reasonable service performance for small-scale HPC applications, thus avoiding the performance penalty usually imposed by the physical network of Clouds \cite{alves_interference-aware_2020}.

        A limiting factor of HPC-Cloud is the abstraction of the underlying Cloud infrastructure from the user, which limits the possibility of tuning the HPC applications. Moreover, while Cloud networks have experienced substantial improvements in network performance (bandwidth/latency), the performance is yet to match HPC performance to optimally host large-scale multi-node workloads. Additionally, having multiple users (multi-tenancy) adds an additional layer of overhead over which the users do not have any control. Cross-interference may happen when applications share a common physical machine, affecting negatively their performances. This problem occurs frequently when high-performance applications are executed in Clouds \cite{alves_interference-aware_2020}. Cloud providers stack several workloads on the same physical resources to explore economies of scale; an approach not always appropriate for HPC applications \cite{netto_hpc_2019}, which raises questions over how to deploy applications properly in an HPC-Cloud converged environment. More research is needed to harmonize HPC and Cloud paradigms “to get the best of both worlds” concerning transparent computing performance, scalability, and resilience \cite{marques_harnessing_2021}.

        The compromise that currently overcomes these issues is the Infrastructure as a Service (IaaS), which is already able to replace small and medium on-premise HPC clusters while providing a lot more flexibility in the number of cores, fluctuating workloads, storage solutions, and other strengths generally associated with the public Cloud \cite{fernandez_evaluation_2022}. Furthermore, simulations show significant improvements in adding elastic jobs with up to 35.34\% higher system utilization, 75.3\% lower runtime, 99.76\% lower wait time, and 75.22\% lower total turnaround time \cite{eberius_evaluating_2023}. Overall, there are aspects that fundamentally clash between HPC and Cloud still today, while others promote convergence with great results.

        \begin{table*}[ht]
\centering
\small
\renewcommand{\arraystretch}{1.5} % Adjust vertical spacing as needed
\resizebox{1\textwidth}{!}{ % Scale the table to 80% of text width, adjust height
\begin{tabular}{l p{5cm} p{5cm} p{5cm}}
    \toprule
     \textbf{Approach} & \textbf{HPC in the Cloud} & \textbf{HPC Plus Cloud} & \textbf{HPC as a Service} \\
    \midrule
    \textbf{Description} & Moving existing HPC applications to cloud environments, leveraging cloud providers' resources and services for execution. \cite{netto_hpc_2019, liu_hpc_2022}  &  Utilizing cloud resources as a supplement to on-premise HPC systems, primarily during peak workloads (Cloud Bursting). \cite{netto_deciding_2015, endo_consideration_2022} & Cloud providers expose HPC resources (hardware, software) as a service, allowing users to access them on demand without managing the underlying infrastructure. \cite{netto_hpc_2018, liu_hpc_2022} \\
    \midrule
    \textbf{Strengths} & \textbf{On-demand access} \cite{netto_hpc_2019}
    \newline \textbf{Scalability} \cite{netto_hpc_2019}
    \newline  \textbf{Cost-effectiveness} \cite{netto_hpc_2019}
    \newline  \textbf{Improved network performance} due to advancements in cloud networking technology. \cite{guidi_10_2021, fernandez_evaluation_2022} & \textbf{Cost-effective} for handling temporary workload spikes. \cite{netto_deciding_2015}
    \newline \textbf{Flexibility} without additional on-premise investments. \cite{netto_deciding_2015}
    \newline \textbf{Seamless integration} with existing HPC workflows. \cite{endo_consideration_2022} & \textbf{On-demand access} to high-performance resources, including specialized hardware. \cite{netto_hpc_2018}
    \newline \textbf{Reduced upfront investment} in hardware and software. \cite{netto_hpc_2018}
    \newline \textbf{Flexibility and scalability}. \cite{netto_hpc_2018}
    \newline \textbf{Access to a wider range of resources} and services. \cite{liu_hpc_2022} \\
    \midrule
    \textbf{Weaknesses} & \textbf{Limited control} over hardware and software configurations. \cite{gupta2014challenges, marathe2013challenges}
    \newline \textbf{Challenges in porting applications} due to compatibility issues. \cite{zhao_resource_2023, liao_bridging_2021, choudhary_exascale_2021}  & \textbf{Potential performance issues} during transitions. \cite{netto_deciding_2015, endo_consideration_2022}
    \newline \textbf{Workflow management challenges} across multiple environments. \cite{balis_scientific_2021}
    \newline \textbf{Requires careful planning and optimization}. \cite{cunha_job_2016} &  \textbf{Cost concerns} for long-running workloads. \cite{netto_hpc_2018, liu_hpc_2022}
    \newline \textbf{Requires specialized knowledge and tools}. \cite{netto_hpc_2018}
    \newline  \textbf{May limit control over resources}. \cite{liu_hpc_2022} \\
    \bottomrule
\end{tabular}
}
\caption{HPC-Cloud Convergence Approaches}
\label{tab:hpc_cloud_approaches}
\end{table*}






    \subsubsection{Hybrid Cloud and HPC Integration: The Best of Both Worlds:} The convergence of HPC and Cloud computing is often realized through hybrid Cloud models, where on-premise HPC systems are integrated with Cloud resources to provide a flexible and scalable computing environment. This approach combines the high performance and dedicated resources of on-premise HPC with the scalability, cost-efficiency, and flexibility of Cloud computing. \\
    \newline
    \textbf{Cloud Bursting: Scaling on Demand:} Cloud bursting allows users to leverage Cloud resources during peak demand periods, supplementing on-premise HPC systems. This approach provides a cost-effective way to handle temporary spikes in workload without investing in additional on-premise infrastructure.
    \begin{itemize}
        \item \textbf{Performance and Cost Analysis:} Studies explore the operational considerations and performance implications of Cloud bursting, evaluating the impact on job execution time and monetary costs using various benchmarks. The results show that Cloud bursting can be an effective strategy for handling peak workloads, but it is important to consider the performance and cost trade-offs \cite{netto_deciding_2015, endo_consideration_2022}.

        \item \textbf{User Transparency:} The importance of user transparency in Cloud-bursting environments is emphasized, recommending subsidizing the cost gap between Cloud and on-premise resources to encourage Cloud usage during peak loads. It also proposes user education and extended job scheduler capabilities to manage job offloading effectively \cite{endo_consideration_2022}.
    \end{itemize}

    \noindent
    \textbf{Workflow Management: Orchestrating Complexity:} The management of complex scientific workflows across hybrid Cloud environments is a critical challenge. These workflows often involve multiple stages, data transfers, and resource dependencies, requiring sophisticated orchestration and management tools.

\begin{itemize}
\item \textbf{Transparent Data Access:} Frameworks for managing workflows across multiple platforms have been proposed, ensuring data consistency and efficient resource allocation. These solutions effectively manage scientific workflows on hybrid Clouds, leveraging Kubernetes for portability and Onedata for transparent data access. A partitioning algorithm balances load distribution and minimizes communication volumes, showing significant potential for optimizing scientific computing tasks across hybrid Cloud environments \cite{balis_scientific_2021}.
\item \textbf{Multiscale Simulations:} The feasibility of using a containerized approach to deploy and execute multiscale scientific workflows in hybrid Cloud environments has been demonstrated. This framework enhances flexibility, reproducibility, and scalability of complex simulations, focusing on an end-to-end modeling case of mayonnaise production. It integrates simulation services at atomistic, mesoscopic, and continuum levels, leveraging containerization for flexible and reproducible deployment of complex multiscale simulations in both on-premise and Cloud systems \cite{elisseev_multiscale_2022}.
\end{itemize}

    \noindent
    \textbf{Software Resource Disaggregation: Leveraging Serverless for HPC:} The integration of serverless computing with HPC environments offers a new approach to resource disaggregation. This approach allows for the dynamic allocation of resources on idle HPC nodes, improving resource utilization and cost efficiency.

    \begin{itemize}
        \item \textbf{FaaS Integration:} Studies demonstrate that serverless computing can dynamically allocate resources on idle HPC nodes, improving resource utilization and cost efficiency. The proposed system leverages Function-as-a-Service (FaaS) to provide fine-grained access to HPC resources, allowing underutilized nodes to be dynamically allocated for short-term tasks without significant hardware modifications \cite{copik_software_2024}.
        \item \textbf{rFaaS:} rFaaS (Remote Function-as-a-Service) is an RDMA-accelerated serverless platform designed for high-performance and latency-sensitive applications. It leverages RDMA to achieve microsecond-scale invocations and efficient data movement, making it suitable for HPC workloads \cite{copik_rfaas_2023}.
        \item \textbf{FuncX:} FuncX is a federated serverless platform that enables the execution of functions across distributed resources. It offers a flexible and scalable solution for HPC workloads, but its invocation latency can be higher due to its distributed nature \cite{chard_funcx_2020}.
        \item \textbf{Wukong:} Wukong is a scalable and locality-enhanced framework for serverless parallel computing. It leverages AWS Lambda and optimizes function placement for low latency and high throughput \cite{carver2020wukong}.
        \item \textbf{Archipelago:} Archipelago is a scalable, low-latency serverless platform that focuses on resource partitioning and decentralized scheduling. It aims to improve the performance and efficiency of serverless applications \cite{singhvi_archipelago_2019}.
        \item \textbf{Performance and Efficiency:}  Another study demonstrates that the system can achieve near-native performance while significantly improving resource efficiency. It highlights substantial cost savings and energy efficiency gains through better resource allocation and reduced idle times \cite{tavares_leveraging_2021}.
        \item \textbf{XaaS:} Aims to provide a unified architecture that leverages the strengths of both HPC and Cloud, enabling users to access high-performance resources transparently, regardless of the underlying cloud or HPC service provider \cite{XaaS-torsten}.
    \end{itemize}

    Overall the work seems to cover extensively the use of serverless computing for reduction of latency and enhancement of scalability for HPC workloads, while achieving near-native performance, Better resource utilization is a byproduct of these efforts, leading subsequently to a cost reduction too. Finally, the development of unified architectures like \textit{XaaS} integrates HPC and Cloud by making HPC resources more accessible and efficient.

    % Imagine harnessing the raw computational muscle of a supercomputer without the hefty maintenance and upfront costs. That's the tantalizing vision of Cloud-HPC convergence, where the boundaries between Cloud and High Performance Computing blur, creating a unified system capable of seamlessly executing both traditional HPC and Cloud applications (Gadban, 2022). This subchapter delves into this transformative journey, exploring its potential, the hurdles it faces, and the exciting research paths leading towards its complete realization.

    % Three distinct paths lead us towards this convergence:

    % HPC in the Cloud: \\
    % Think of this as migrating existing HPC applications to the Cloud's vast resource pool. Imagine simulating complex weather patterns or analyzing massive datasets through Cloud-based computing power, previously accessible only through dedicated HPC clusters.
    % HPC Plus Cloud: \\
    % This approach acts as a sidekick, supplementing on-premise HPC resources with Cloud infrastructure during peak demand periods. Think of it as Cloud bursting, where intensive workloads temporarily spill over to the Cloud, ensuring smooth operations without exceeding on-premise capacity.
    % HPC as a Service: \\
    % Here, the Cloud becomes the primary provider, offering HPC resources directly through Cloud services. Imagine accessing specialized hardware like GPUs or high-performance storage on demand, eliminating the need for upfront investment and ongoing maintenance.

    % Each approach offers unique advantages in terms of resource allocation and user experience. For instance, containerization technologies offer lightweight virtualization, streamlining application deployment and management in the Cloud \cite{marques_harnessing_2021}. This flexibility empowers researchers and developers to experiment with diverse hardware configurations, fostering innovation without breaking the bank.

    % However, the path to seamless convergence isn't without its roadblocks. The Cloud's inherent abstraction layers, designed for general-purpose computing, can sometimes limit the fine-tuning capabilities crucial for optimizing HPC applications, potentially impacting performance \cite{netto_hpc_2019}. Additionally, the high-bandwidth demands of tightly coupled HPC workloads can outpace existing Cloud network capabilities, leading to bottlenecks and performance degradation \cite{alves_interference-aware_2020}.

    % But the research community is undaunted, actively tackling these challenges. One promising solution is **elasticity in HPC**, enabling applications to dynamically adjust resource usage based on workload demands \cite{eberius_evaluating_2023}. Imagine an HPC application automatically scaling up resources during intensive calculations and scaling down during idle periods, optimizing both performance and cost efficiency.

    % Another exciting avenue is \textbf{resource disaggregation}, where individual resources like CPU, memory, and storage are decoupled and allocated independently within the Cloud infrastructure \cite{michelogiannakis_case_2022}. This granular control allows for more efficient utilization of specialized hardware, catering to the diverse needs of various HPC applications.

    % Furthermore, researchers are developing dedicated \textbf{high-performance Cloud networks} specifically designed to handle the intense data movement demands of HPC workloads. These networks promise to eliminate bandwidth bottlenecks, unlocking the full potential of Cloud-based HPC for large-scale simulations and data analysis.

    % The impact of Cloud-HPC convergence extends far beyond technical considerations. Imagine democratizing access to cutting-edge computing power, enabling smaller research institutions and startups to tackle complex problems previously restricted to large organizations with dedicated HPC infrastructure. This democratization opens doors to groundbreaking discoveries across various fields, from drug development and materials science to climate modeling and artificial intelligence.

    % The future of Cloud-HPC convergence is brimming with possibilities. With dedicated research efforts addressing existing challenges and unlocking new capabilities, this transformative technology holds the potential to revolutionize how we approach computational challenges, ushering in an era of faster scientific breakthroughs, groundbreaking innovations, and a more collaborative and accessible research landscape.

    % This expanded version incorporates specific examples, highlights the broader societal impact of Cloud-HPC convergence, and offers a more holistic perspective on its future potential.





    % ################################################################################
\subsection{Reproducibility and Transparency}
\label{sec:reproducetransparency}
    The portability of HPC applications to Cloud environments is a significant challenge. HPC systems often rely on highly optimized, vendor-specific tools and libraries, which can be challenging to port to Cloud platforms. This incompatibility can lead to performance degradation and necessitate significant code modifications \cite{zhao_resource_2023, liao_bridging_2021, choudhary_exascale_2021}.

    Containerization technologies, such as \textit{Docker} and \textit{Singularity/Apptainer}, offer a solution for portability. By encapsulating applications and their dependencies within isolated environments, containers can be deployed across different platforms with minimal effort. However, containerization itself introduces overhead, and the performance of containers on HPC systems can vary significantly depending on the container engine and the underlying hardware \cite{zhang_performance_2021, zhang_framework_2023, zhao_resource_2023}.
    Efforts to standardize container formats and APIs, such as the Open Container Initiative (OCI), are crucial for improving portability. However, these standards do not yet address all aspects of HPC-specific requirements, such as the need for efficient access to specialized hardware accelerators \cite{zhao_resource_2023}.

    Ease of use is a major factor in enabling portability. Complex setup procedures are a barrier for new users and limit the adoption of HPC resources. Cloud platforms, on the other hand, are generally designed to be user-friendly, with intuitive interfaces and automated tools. This accessibility makes Cloud computing appealing to a broader range of users [1, 3, 10, 15, 17, 22]. Additionally, a set of methodological principles exists that attempt to further foster reproducibility in cloud computing research, which includes repeated experiments, workload and configuration coverage, experimental setup description, open access artifact, probabilistic result description of measured performance, statistical evaluation, reporting of measurement units and cost model \cite{iosup-reproducibility}. Ideally, a converged system should offer a consistent and intuitive interface that hides the underlying complexity of both HPC and Cloud infrastructure.




    % ################################################################################
\subsection{Economic and Market Dynamics}
\label{sec:economicmarketdynamics}
        % “Most research efforts in HPC-Cloud aim to understand the cost-benefit of moving resource-intensive applications from on-premise environments to public Cloud platforms. Industry trends show hybrid environments are the natural path to get the best of the on-premise and Cloud resources—steady (and sensitive) workloads can run on on-premise resources and peak demand can leverage remote resources in a pay-as-you-go manner.” ([Netto et al., 2019, p. 81](zotero://select/library/items/PDHAK8LB)) ([pdf](zotero://open-pdf/library/items/43NTJN9N?page=1))


    The majority of research in the field of HPC Cloud is aimed at understanding the cost-benefit analysis of migrating resource-intensive applications from on-premise environments to public Cloud platforms \cite{netto_hpc_2019}. These environments are managed by batch schedulers that queue user requests to run jobs when resources are heavily used. However, without Service Level Agreements (SLAs), users have no visibility or concerns about the costs of running jobs, leading to potential resource wastage and poor quality of service \cite{netto_hpc_2019}.


     Steady and sensitive workloads can run on on-premise resources, while peak demand can leverage remote resources in a pay-as-you-go manner \cite{netto_hpc_2019}. However, Cloud providers often offer machines with a higher number of CPU threads and more memory than necessary, increasing the costs of the cluster \cite{gurgul_cloud-native_2023}. In contrast, GPU-based clusters with thousands of threads per node may seem superior, but they also have their limitations \cite{gurgul_cloud-native_2023}. The predictability of running computations in the Cloud is lower than in on-premise clusters due to the multi-tenant nature of Cloud and the potential for many concurrent workloads scheduled on the same hardware \cite{gurgul_cloud-native_2023}. One solution to this problem is the exclusive rental of machines in the Cloud (bare-metal). However, this service is currently limited in its availability and comes at a considerably higher price compared to regular nodes \cite{gurgul_cloud-native_2023}. Furthermore, the cost of running workloads and storing their results on-premises and in the Cloud can be prohibitive when the cluster is heavily used, especially when considering other expenses like traffic and storage \cite{gadban_analyzing_2022}. Despite these challenges, many organizations have been moving their applications or even their entire data centers to the Cloud, motivated by the economic advantages in terms of real estate, equipment costs, and manpower compared to private data centers \cite{alashaikh_survey_2021}. This shift has been facilitated by the high degree of flexibility offered by Cloud data centers, where services and applications can be deployed, terminated, migrated, and replicated on demand \cite{alashaikh_survey_2021}. However, it has been shown that users, not just service providers, actively affect the environmental impact of such systems. This highlights the need for improved mechanisms that enable interaction between providers and users to optimize resource usage, specifically targeting HPC due to additional constraints that traditional data centers do not experience \cite{silva_review_2024}. The increasing usage rate in HPC as a Service (HPCaaS) can contribute to more sustainable HPC. This can be achieved not only by taking advantage of HPC-Cloud sites with better conditions for free cooling and availability of Renewable Energy Sources (RES) but also by using \textit{spatio-temporal workload shifting} with regional grid carbon intensity as a signal, similar to what is already being done in traditional Cloud Computing \cite{silva_review_2024}.

     The HPC-Cloud convergence, currently, is limited by the fact that HPC communities are often conservative in adopting new technologies and alternative paradigms, which is rooted in the high cost of HPC systems and the need to ensure that new technologies do not compromise performance. The cost-effectiveness of Cloud-based HPC solutions is a major driver for convergence. Cloud computing offers a pay-as-you-go model, allowing users to access high-performance resources on demand without the upfront investment required for traditional HPC systems. This flexibility, combined with the potential for resource optimization, makes Cloud-based HPC an attractive option for researchers and organizations with varying computational needs. \\


    \subsubsection{Cost Efficiency and Resource Management: Optimizing for Value}
    \label{sec:costefficiency}



    \paragraph{Instance Selection:} Studies like \cite{walkup_best_2022} provide guidance on selecting the most cost-effective VM instances for different HPC workloads. The choice between single large instances and multiple smaller instances depends on the specific application requirements and communication patterns.
    \begin{itemize}
    \item{Large Shared-Memory Instances} For latency-sensitive applications, large shared-memory instances outperform multiple smaller instances due to reduced communication overhead. This is because data can be accessed directly within the VM's memory, eliminating the need for network communication between different instances.
    \item{Cost-Effectiveness} For configurations with the same total number of vCPUs and memory, the cost-effectiveness is similar. This means that users can choose the configuration that best suits their specific needs without significant cost implications.
    \end{itemize}
    \paragraph{Resource Disaggregation:}
    \begin{itemize}
        \item {Intra-Rack Disaggregation} Studies like \cite{michelogiannakis_case_2022} demonstrate that disaggregation can significantly reduce resource usage and improve cost efficiency. The study analyzes the utilization of memory and NIC resources in an HPC environment, finding that a significant proportion of time memory occupancy is well below capacity. The research quantifies the benefits of intra-rack resource disaggregation, particularly focusing on reducing memory and NIC resources.
        \item {Photonics Integration} \cite{michelogiannakis_efficient_2023} explores the use of co-packaged dense wavelength division multiplexing (DWDM) photonics for intra-rack resource disaggregation. This approach provides a flexible and high-bandwidth solution to meet the stringent requirements of HPC applications, achieving an average application speedup of 11\% (up to 46\%) for CPU benchmarks and 61\% for GPU benchmarks compared to systems using electronic switches.
    \end{itemize}

    \paragraph{Automated Optimization - Intelligent Resource Management:}

    Several studies propose automated systems for detecting and rectifying inefficiencies in Cloud resource configurations. These systems leverage machine learning and data analysis techniques to identify suboptimal configurations and suggest improvements, leading to significant cost savings and performance enhancements \cite{greneche_methodology_2022, michalewicz_attempt_2023,giannoutakis_evaluation_2021}.
    \begin{itemize}
        \item Automated Detection and Rectification: highlights the potential for significant cost savings through automated optimization of Cloud resource usage. The study identifies common pitfalls in resource configuration, proposes automated detection and rectification techniques, and demonstrates significant cost savings and performance improvements through real-world case studies.
        \item Performance and Scalability: A study shows that the automated detection and rectification system can scale across different Cloud environments and HPC workloads, maintaining effectiveness in diverse scenarios. The system achieves substantial performance gains by aligning resource allocation with workload requirements, resulting in faster job completion times and improved resource utilization.
    \end{itemize}

    Moving to the Cloud involves a shift from capital expenditure (CapEx) to operational expenditure (OpEx). Capital expenditure (CapEx) involves large, upfront investments in physical assets like equipment or infrastructure, which are depreciated over time. Operational expenditure (OpEx) refers to ongoing costs for running a business, such as utilities and service subscriptions. In cloud computing, OpEx represents the shift from purchasing and maintaining hardware (CapEx) to paying for cloud services on a flexible, usage-based model, reducing the need for significant upfront spending. This transition is attractive for organizations with limited capital \cite{netto_hpc_2019,varghese_next_2018,gill_taxonomy_2018,ahn_flux_2020,rodriguez_container-based_2019,netto_hpc_2019,gadban_analyzing_2021,gadban_investigating_2020}. However, the OpEx model can be more expensive in the long run, especially for compute-intensive workloads.
    Migrating HPC workloads to the Cloud can involve significant costs for software and hardware adaptation. This includes the cost of porting applications, acquiring new licenses, and potentially purchasing new hardware. The transition to a converged environment also requires training for both users and administrators. This can be a significant cost factor, especially for organizations with limited resources. Overall, the costs of the transition to a converged system are unclear at this moment.









% \begin{table*}[h!]
% \centering
% \small
% \caption{Summary of Key Points for Slides on HPC-Cloud Convergence}
% \begin{tabular}{p{3.5cm}|p{10.5cm}}
% \toprule
% \textbf{Slide Topic} & \textbf{Key Points} \\
% \midrule
% \textbf{Technological Aspects and Infrastructure} &
% \makecell[l]{\textbf{Closing the Performance Gap:} Hardware Advancements, Network Performance (InfiniBand, RDMA, CXL), Specialized Hardware (GPUs, TPUs, FPGAs), Bare-Metal Servers \\
% \textbf{Enabling Portability and Flexibility:} Containerization (Docker, Apptainer), Elasticity (DMR API, PMIx), Cloud Storage for HPC (S3, Azure Blob) \\
% \textbf{Optimizing Resource Utilization:} Resource Disaggregation, Improved rack utilization, Challenges in large-scale implementation} \\
% \midrule
% \textbf{Performance, Configuration, and Deployment Models} &
% \makecell[l]{\textbf{Approaches to HPC-Cloud Convergence:} HPC in the Cloud, HPC plus Cloud, HPC as a Service, Mixed Systems \\
% \textbf{Challenges and Limitations:} Abstraction limits tuning, Network limitations, Multi-tenancy issues, Research needs \\
% \textbf{IaaS:} Replacing on-premise clusters, Flexible core count, Elastic jobs with public Cloud benefits} \\
% \midrule
% \textbf{Hybrid Cloud and HPC Integration} &
% \makecell[l]{\textbf{Hybrid Cloud:} Combines on-premise HPC with Cloud, Cloud Bursting, Workflow Management \\
% \textbf{Serverless Computing:} Dynamic resource allocation, Improved utilization, FaaS Integration \\
% \textbf{Towards Unified Architecture:} XaaS with performance-portable containers, Flexibility with serverless computing} \\
% \midrule
% \textbf{Reproducibility and Transparency} &
% \makecell[l]{\textbf{Portability Challenges:} Vendor-specific tools, Performance degradation, Code modifications \\
% \textbf{Containerization:} Docker, Singularity/Apptainer for portability, Overhead and performance considerations \\
% \textbf{Standardization:} Open Container Initiative (OCI), Hardware accelerator access issues} \\
% \midrule
% \textbf{Economic and Market Dynamics} &
% \makecell[l]{\textbf{Cost-Benefit Analysis:} Migrating HPC to Cloud, Resource wastage, Hybrid environments \\
% \textbf{Cloud Provider Considerations:} Cost optimization, Oversized machines, Multi-tenancy impact, Bare-metal rental costs \\
% \textbf{Shift to Cloud:} Economic advantages, Flexibility, Environmental impact optimization} \\
% \bottomrule
% \end{tabular}
% \end{table*}





    % Cost-benefit analysis of moving resource-intensive applications from on-premise environments to public Cloud platforms.
    % Hybrid environments as the natural path forward.
    % Challenges in resource allocation and cost management.
    % Impact of multi-tenancy on cost predictability.
    % Exclusive rental of machines as a solution to cost predictability.
    % Cost of running workloads and storing results on-premises and in the Cloud.
    % Economic advantages of moving applications or data centers to the Cloud.
    % Flexibility offered by Cloud data centers.
    % Users' impact on the environmental sustainability of Cloud systems.
    % Mechanisms for optimizing resource usage and sustainability.
    % Spatio-temporal workload shifting for sustainability.
    % Future trends in HPC systems featuring GPUs, accelerators, and new memory technologies.

\begin{figure*}[]
      \centering
      \resizebox{0.9\textwidth}{!}{ \includesvg{tex files/tables/timeline.drawio.svg} }
      \caption{Timeline of HPC-Cloud Convergence}
      \label{fig:hpc_cloud_timeline}
    \end{figure*}

\section{Discussion}

% \subsection{Timeline of cloud and HPC adoption}
\subsection{Emerging trends}
\label{subsec:emerging_trends}
    % The intersection of HPC and cloud computing is rapidly evolving, driven by a wave of innovative trends that promise to revolutionize how scientific research, engineering design, and data-intensive applications are conducted. This convergence is fostering a new era of flexibility, scalability, and efficiency in HPC environments. \autoref{fig:hpc_cloud_timeline} illustrates the approximate estimation of the periods in which concrete trends were identified in the literature. After the topics were identified in the collected literature, an additional search iteration was performed on scientific databases (mostly Scopus). This iteration was meant to observe when a significant growth in publications for the specific topic was starting to gain popularity (\autoref{fig:enter-label}). There are some areas that could be identified with a temporal offset between the emerging trend and the actual first attempts. An example is about the identified growth in topic of Serverless computing, where some influential works have been done in 2017, while the topic growth has been labeled in 2021 \cite{van2017spec}. Thus, the timeline estimation is not exact, but it gives a clear overview of the progress in the convergence process.


    The intersection of HPC and cloud computing is rapidly evolving, driven by innovative trends that are set to revolutionize scientific research, engineering design, and data-intensive applications. This convergence is ushering in a new era of flexibility, scalability, and efficiency within HPC environments. Figure \autoref{fig:hpc_cloud_timeline} provides an approximate timeline of when these trends were first identified in the literature. After initial identification, an additional search was conducted in scientific databases (primarily Scopus) to determine when significant growth in publications on these topics began to gain momentum (\autoref{fig:enter-label}). Some trends show a temporal gap between their emergence and the first attempts at implementation. For instance, serverless computing saw influential works in 2017, but significant growth in the topic wasn't observed until 2021 \cite{van2017spec}. While the timeline is not exact, it offers a clear overview of the progress in this convergence.

    \begin{figure}[H]
        \centering
        \includesvg[width=0.5\linewidth]{misc/time_process.svg}
        \caption{Estimating emergence of a topic trend}
        \label{fig:enter-label}
    \end{figure}



    \subsubsection{Resource Disaggregation: Optimizing Hardware Utilization}
    A key trend is resource disaggregation, which breaks down traditional HPC node configurations into smaller, specialized units, maximizing hardware utilization and enhancing flexibility. This approach allows for optimized resource allocation based on specific workload requirements, reducing overheads associated with underutilized components. Intra-rack disaggregation \cite{michelogiannakis_case_2022} and software resource disaggregation \cite{copik_software_2024} are key examples of this trend, enabling HPC systems to adapt dynamically to diverse and evolving demands.
    \subsubsection{Software-Defined Infrastructure: Dynamic and Flexible HPC Environments}
    Software-defined infrastructure (SDI) is transforming HPC environments by creating dynamic, flexible, and multi-tenant clusters that seamlessly integrate HPC and cloud technologies. Initiatives like vClusters \cite{alam_versatile_2023} leverage Infrastructure-as-Code (IaC) to automate resource provisioning and configuration, enabling DevOps practices and enhancing agility in HPC resource management.
    \subsubsection{HPC-Cloud Storage Convergence: Efficient Data Management}
    The convergence of HPC and cloud storage is revolutionizing how data-intensive scientific workflows are managed. This integration minimizes data transfer overhead, expands storage scalability, and provides cost-effective data management solutions. Libraries like S3embedded \cite{gadban_analyzing_2022} are paving the way for seamless integration of cloud storage with HPC systems, enhancing the efficiency of data-intensive scientific applications.
    \subsubsection{Edge and Serverless Computing: Expanding HPC Capabilities}
    \label{subsec:edge_computing}
    The emergence of serverless computing platforms like AWS Lambda and OpenWhisk has opened new possibilities for running HPC workloads in the cloud. These platforms offer a pay-per-use model, allowing users to access computing resources on demand without managing infrastructure. Studies like \cite{chahal_isesa_2022} and \cite{copik_rfaas_2023} demonstrate that serverless platforms can achieve significant performance improvements for specific HPC applications, particularly those with short execution times and high parallelism. The paper introduces rFaaS, an RDMA-accelerated Function-as-a-Service (FaaS) platform designed to enhance serverless computing for high-performance and latency-sensitive applications. By integrating RDMA and implementing new resource management policies, rFaaS achieves microsecond invocation latencies and improves overall performance, making serverless computing viable for HPC workloads. \cite{chahal_isesa_2022} additionally introduces the intelligent Serverless Scalable Architecture (iSeSA) framework, which facilitates the migration of HPC and Artificial Intelligence (AI) workloads to serverless platforms. The framework leverages machine learning (ML) and deep learning (DL) techniques to optimize workload execution on serverless platforms, providing scalability and flexibility. The iSeSA framework automatically decides the execution environment (on-premise or serverless) based on resource demand, cost, and service level agreements (SLAs).

    \subsubsection{AI and Machine Learning: Intelligent Optimization}
    \label{subsec:ai_ml}
    The integration of AI and machine learning into HPC-Cloud systems is a significant emerging trend, with applications ranging from optimizing resource allocation to powering new scientific applications.  AI is being used to enhance the performance and efficiency of HPC-Cloud systems in various ways. The Intelligent Serverless Scalable Architecture (iSeSA) framework \cite{chahal_isesa_2022} demonstrates how AI can be used to predict workload execution times and costs, enabling more efficient allocation of serverless resources for HPC and AI workloads. This framework also showcases the use of AI models for accurate temperature predictions for predictive maintenance, demonstrating the suitability of serverless platforms for AI inference tasks \cite{chahal_isesa_2022}. Another promising application is the use of machine learning to identify and rectify inefficiencies in cloud resource configurations for HPC, leading to significant cost savings and performance improvements \cite{michalewicz_attempt_2023}.
    Research also indicates the potential for AI-driven resource-aware scheduling algorithms \cite{silva_review_2024}, enabling dynamic allocation of resources based on real-time workload demands in distributed HPC environments. AI can also be used to predict the performance of HPC applications on different cloud instances and automatically place workloads on the most suitable resources, significantly improving performance and reducing costs \cite{liang_scalable_2022}. Dynamic resource scaling, where AI adjusts the number of resources allocated to an HPC application based on real-time performance metrics, can help to ensure optimal performance while minimizing costs \cite{eberius_evaluating_2023}. AI can also be used to optimize network traffic flow and reduce latency in cloud environments, which is particularly important for HPC applications that require high-bandwidth communication \cite{alves_interference-aware_2020}. Furthermore, AI can be used to detect and recover from failures in HPC-Cloud systems, ensuring the reliability and availability of HPC applications. Finally, AI can be used to optimize the energy consumption of HPC-Cloud systems, reducing costs and minimizing the environmental impact of HPC \cite{silva_review_2024}.
    Examples of AI-powered HPC applications are emerging in diverse fields. The BRICS CloudHPC project showcases a cerebrovascular segmentation algorithm accelerated by GPUs, demonstrating the potential of AI-powered applications in HPC-cloud environments \cite{noauthor_harnessing_nodate}.
    Looking forward, the field is recognizing the need for improved Uncertainty Quantification (UQ) methods to better manage uncertainties in cloud environments \cite{kabir_uncertainty-aware_2022}. AI and ML techniques are seen as promising solutions for this challenge. The "HPC Forecast: Cloudy and Uncertain" article \cite{reed_hpc_2023} further highlights the increasing importance of AI technologies in shaping the future of HPC. The survey emphasizes the development and application of AI and machine learning for predicting and managing uncertainties in cloud computing \cite{kabir_uncertainty-aware_2022}. This includes areas like intelligent cloud management, where AI can be used to predict and adapt to changing resource availability and workload demands.




\subsection{HPC and Cloud Convergence: A New Paradigm}
\label{subsec:convergence}

    As mentioned in \autoref{sec:perfconfiguration}, there seem to be three main options in HPC-Cloud: HPC in the Cloud, HPC Plus Cloud, and HPC as a Service. Here we discuss why HPC Plus Cloud or a converged option might be more future-proof. This approach recognizes the inherent strengths of both environments and seeks to leverage them for optimal results.
    On-premise HPC offers dedicated resources, high-bandwidth interconnects, and specialized hardware optimized for demanding workloads. Cloud computing, on the other hand, provides scalability, flexibility, and pay-per-use models, ideal for handling fluctuating demands and resource-intensive tasks. Convergence allows users to harness the best of both worlds, creating a hybrid solution that balances performance, cost-effectiveness, and flexibility.
    A complete cloud transition presents several challenges. The cost of running large-scale HPC applications on the cloud can be prohibitive, especially for long-term projects. Concerns around data security and privacy can arise when moving sensitive HPC data to public cloud environments. Additionally, complete cloud migration might limit control over the infrastructure and software environment, potentially disrupting existing workflows and research practices.
    Convergence allows for a gradual adoption of cloud services while leveraging existing HPC investments. This approach facilitates the creation of hybrid workflows, where on-premise resources are combined with cloud services for specific tasks or workloads. Cloud bursting allows scaling up to cloud resources during peak demand periods, providing flexibility without the commitment of a total cloud transition. Convergence also enables incremental integration, where cloud-based storage solutions, tools, and services are incorporated within existing HPC environments.
    Several entries in the literature highlight the benefits of this hybrid approach. For example, \cite{balis_scientific_2021} demonstrates the effectiveness of managing complex scientific workflows across hybrid cloud environments, leveraging cloud bursting for scalability and flexibility. The research exploring a hybrid cloud-HPC architecture for seamless job execution, \cite{endo_consideration_2022} enables users to leverage cloud resources for peak demands without altering their workflows. The vClusters project integrates HPC and cloud technologies, creating a versatile, multi-tenant cluster on the Alps supercomputer, showcasing the benefits of a hybrid approach \cite{alam_versatile_2023}.
    There are more examples of literature, directly or indirectly, suggesting that the hybrid modality has the most potential to satisfy cost, usability, portability, transition, and performance criteria.



Some of the observations found in the literature include the following.

\begin{itemize}
    \item Cloud services can enhance existing HPC systems by handling specific tasks within a larger infrastructure, demonstrating the potential for cloud computing to complement and extend HPC capabilities \cite{colonnelli_bringing_nodate}.
    \item There is a need for specialized tools to predict how well HPC applications will perform on cloud platforms, suggesting that hybrid systems might require different approaches than running solely on the cloud \cite{noauthor_hpccloud_nodate}.
    \item Choosing the best place to run an HPC job—on-premise or in the cloud—is complex, and hybrid environments require careful planning to ensure optimal performance \cite{cunha_job_2016}.
    \item Using a combination of cloud and HPC resources can be beneficial for managing large, complex scientific workflows, offering flexibility and the ability to scale up quickly \cite{balis_scientific_2021}.
    \item Cloud technology can make HPC systems more efficient and flexible, supporting the idea of combining these two technologies for better performance \cite{noauthor_harnessing_nodate}.
    \item HPC and cloud systems can be integrated to create a seamlessly functioning hybrid environment \cite{noauthor_consideration_nodate}.
    \item Combining cloud storage and HPC can improve the handling of massive amounts of data in research \cite{gadban_analyzing_2022}.
\end{itemize}

In conclusion, the literature emphasizes that the convergence of HPC and cloud computing presents a more practical and beneficial approach compared to a complete cloud transition. The trend also shows that in the span of the last decade, the challenges seem to be rooted in finding correct configurations to ensure maximum performance. However, the challenges are being solved, and the mixed nature of the two converging paradigms allows for a broad spectrum of innovations. The hybrid approach allows researchers and developers to harness the best of both worlds, achieving optimal performance, cost-effectiveness, and flexibility for diverse HPC workloads.



\subsection{A Glimpse into the Future of Data-Intensive Science: The Power of Convergence}
\label{subsec:future}

    The convergence of HPC and cloud computing, as explored in this survey, has the potential to transform the landscape of HPC, and it is doing so already. Here’s a vision of this future:\\
    \newline
    \newline
    \noindent
    \textit{Accelerated Scientific Discovery}: A converged system will empower scientists to tackle larger, more complex scientific problems, leading to faster and more impactful discoveries. The ability to access immense computational and storage resources on demand, coupled with efficient data sharing, will enable simulations of unprecedented scale and complexity \cite{gadban_analyzing_2022, netto_hpc_2019}.
    \newline
    \newline
    \noindent
    \textit{Breakthroughs in Emerging Fields}: This convergence will fuel innovation in areas such as artificial intelligence, machine learning, and big data analytics. The access to massive computational resources and the capability to seamlessly share and analyze massive datasets will accelerate research and development in these fields \cite{silva_review_2024, liang_scalable_2022}.
    \newline
    \newline
    \noindent
    \textit{Democratization of HPC}: Converged systems will democratize access to HPC resources, allowing researchers and organizations with limited budgets to participate in cutting-edge research and development efforts. The ability to scale resources dynamically based on demand will make HPC more accessible to a broader community \cite{reed_hpc_2023}.
    \newline
    \newline
    \noindent
    \textit{Greater Efficiency and Cost Reduction}: Converged systems can optimize resource utilization by leveraging the elasticity of cloud computing and the high performance of HPC platforms. It will reduce the cost of running scientific workflows, improve system utilization, and lower energy consumption. This will enable organizations to allocate resources more effectively and achieve significant cost savings. The journey towards a truly converged HPC-Cloud ecosystem will involve overcoming numerous technical and organizational challenges. However, by leveraging the insights and solutions proposed in this dissertation, we can foster a new era of scientific exploration and unlock the full potential of data-intensive science \cite{noauthor_hpccloud_nodate, alam_versatile_2023, gadban_analyzing_2022, liang_scalable_2022, netto_deciding_2015}. \\

    What will contribute to the convergence of HPC and cloud computing? The following factors will play a crucial role in shaping the future of scientific computing:
    \newline
    \newline
    \noindent
    \textit{The Role of Serverless Computing}: Serverless computing platforms like AWS Lambda offer a highly scalable and cost-effective solution for HPC workloads, especially those with dynamic and unpredictable resource demands \cite{copik_rfaas_2023, petrosyan_serverless_2022, chahal_isesa_2022}.
    \newline
    \newline
    \noindent
    \textit{Edge Computing}: Edge computing can further enhance HPC-Cloud convergence by bringing computation closer to the data source and reducing network latency and bandwidth requirements \cite{greneche_methodology_2022, kuity_investigating_2023, guo_understanding_2023, endo_consideration_2022}.
     \newline
     \newline
    \noindent
     \textit{Quantum Computing}: Quantum computing is poised to revolutionize HPC, and cloud providers are already offering access to quantum computers and simulators. This opens up new possibilities for solving complex scientific problems that are intractable for classical computers \cite{michalewicz_attempt_2023, ruhela_characterizing_2021}.

\section{Litmitations and Future Work}
\label{sec:future}

This review only touches on a fraction of the material available on the convergence of HPC and Cloud, constrained by the project's scope and time. The primary challenge lies in the vastness of the field, making it difficult to encapsulate within a short report. Each macro-category covered could easily warrant a dedicated literature survey. Therefore, this review serves as a high-level guide to key trends and topics. While we aimed to include significant aspects within these macro-categories, the following suggestions outline potential directions for future research.
\newline
\newline
\noindent
\textit{Security and Trust}: Moving to a converged environment necessitates robust security measures to protect sensitive data and ensure the integrity of the data ecosystem. It requires thorough analysis of risks associated with cloud deployment, including external threats, supply chain attacks, internal threats, and misconfiguration vulnerabilities. The challenge lies in balancing security with the need for efficient and scalable data access.
\newline
\newline
\noindent
\textit{Data Management and Access}: Maintaining a consistent and reliable data environment in a converged system requires sophisticated management tools and techniques. Efficient data management encompasses data replication, data redundancy, data integrity, and data recovery. A comprehensive set of tools and
procedures needs to be defined to address these challenges, specifically for handling the data lifecycle, which includes data creation, data use, data archival, and eventually data deletion.
\newline
\newline
\noindent
\textit{Resource Allocation and Scheduling}: Orchestrating the execution of scientific workflows in a converged HPC-Cloud environment demands intelligent resource allocation and job scheduling mechanisms. It requires the ability to dynamically allocate resources based on job characteristics, prioritize jobs, optimize performance, minimize costs, and effectively handle system heterogeneity. This necessitates the development of hybrid resource managers capable of adapting to dynamic resource availability patterns and the diverse needs of HPC and cloud applications.
\newline
\newline
\noindent
\textit{Application Adaptability and Portability}: HPC applications often need to be adapted to run seamlessly in a converged environment. While the use of wrapper libraries like S3Embedded simplifies the transition, many applications rely on specific storage interfaces and protocols. It is crucial to develop techniques and tools that enable seamless data access and portability for HPC applications across different platforms. \\

\section{Conclusion}
% \todo{Conclusion of your investigation, including (i) few sentence to introduce the topic considered, (ii) short summary of the research method followed, (iii) main findings, and (iv) potential implications and future work}
The convergence of HPC and Cloud computing presents a compelling vision for the future of scientific computing. It has the potential to transform data-intensive science, enabling accelerated discoveries, fostering innovation, and democratizing access to high-performance computing resources.
The journey towards this converged future is just beginning. The insights and solutions presented in this dissertation offer a roadmap for navigating this exciting landscape and building a future where HPC and cloud computing work in harmony to advance the frontiers of scientific knowledge.


\bibliographystyle{ACM-Reference-Format}
\bibliography{references2}
% \nocite{*}

%
% If your work has an appendix, this is the place to put it.
%\appendix


\end{document}
